% !TeX root = ../main.tex
\section{Introduction}
\label{sec:intro}

Motivational Interviewing \citep[MI;][]{miller1991mi} is a counselling style for helping
people resolve ambivalence about change. It is defined less by what the counsellor says
than by how they say it: open questions, reflective listening, affirmation of the client's
own strengths, and a studied refusal to argue for change on the client's behalf. That
makes MI an attractive and unusually demanding target for goal-oriented dialogue systems.
The behaviour to be learned is multi-turn, the reward is delayed, and expert-annotated
data is scarce and expensive.

The now-standard workaround is to close the loop with language models: a large model plays
the patient, another large model plays an evaluator that fills in validated MI
questionnaires, and a small model is trained against the resulting score
\citep{yosef2024assessing, baruch2025pto}. This is attractive because it is cheap and
because the questionnaires are instruments clinical researchers already trust. It is also
a reward-model design in which the reward model is a black box that was never trained to
be robust to optimization pressure --- and the training loop pushes on it for ten
iterations.

This paper asks two questions about that loop, and answers them under tightly matched
conditions. First: \textbf{how should the reward be turned into a policy update?} We compare
Preference-Tree Optimization (PTO), which grows a \emph{sequential} search over therapist
turns and trains with DPO \citep{rafailov2023dpo} on the surviving preference pairs,
against Group-Relative Policy Optimization \citep[GRPO;][]{shao2024deepseekmath}, which
samples an i.i.d.\ group of completions per prompt and trains on standardized advantages.
Second: \textbf{what do the resulting gains mean?} A questionnaire score that goes up is a
claim about the therapist only to the extent that the questionnaire-filler is not being
gamed.

We hold everything except the two answers fixed: the same Llama-3.2-1B base policy, the
same simulated patients, the same \primaryjudge{} oracle and reward composition, the same
per-prompt candidate budget ($M = G = 8$), the same minimum-context filter, and the same
look-ahead depth ($K=0$). Both arms run ten iterations, regenerating their training data
from the current policy each time; each of the 22 resulting model states is evaluated on
96 patient personas and scored on eight instruments --- five global rubrics, a patient
change-talk coder, an MI-inconsistency coder, and the MITI technique ratios.

\paragraph{Contributions.}
\begin{enumerate}\itemsep2pt
  \item \textbf{A matched comparison of preference-tree and group-relative optimization
  for MI.} Both methods produce large gains over the base policy; PTO leads at the matched
  ten-iteration endpoint (\QQ{} $4.26$ vs $3.75$; paired $+0.51$, $\dz{}=0.73$,
  Holm $p<0.001$) and climbs monotonically, whereas GRPO peaks at iteration~8 and regresses
  (\S\ref{sec:results}).

  \item \textbf{A characterisation of what the reward actually purchases.} In both arms the
  therapist drifts toward affirmation and away from inquiry --- affirmations per turn rise
  ${\sim}5\times$ while literal question rates fall --- and MI-inconsistency rises with the
  scores it is supposed to contradict. We trace this to specific items of the reward
  questionnaire, not only to the optimizer (\S\ref{sec:hacking}).

  \item \textbf{Gain retention under a decoupled judge as a reward-hacking test.} We
  re-score every conversation with a judge from a different model family that never played
  the patient and never touched training. Contrast signs survive ($88.3\%$ overall,
  $98.9\%$ at $|\Delta|\!\ge\!0.50$), but the \emph{share of each arm's gain the held-out
  judge still credits} does not: it is an onset curve that separates the arms from
  iteration~3 and ends at $0.80$ (PTO) versus $0.28$ (GRPO). This is a cheap, general
  diagnostic for any LLM-judge-trained system (\S\ref{sec:validity}).

  \item \textbf{Evidence that the method gap is about exploration, not the loss.} Putting
  both methods' updates on one probe, swapping the weighting rule on the \emph{same}
  candidate groups barely moves the update direction, while holding the rule fixed across
  each method's \emph{own} groups leaves them as far apart as when trained
  (\S\ref{sec:mechanism}; full construction and per-iteration series in
  Appendix~\ref{app:mechanism}).
\end{enumerate}

The practical reading for the CLPsych community is deflationary and, we think, useful: with
an LLM oracle in the loop, ``all rubrics went up'' is not evidence of clinical skill, and
the difference between two RL algorithms may be a difference in what conversations they
bother to look at.
